% LaTeX Curriculum Vitae -- research / industry version
%
% Based on a template by Jason Blevins <jrblevin@sdf.lonestar.org>
% http://jblevins.org/projects/cv-template/
%
% This is the version targeted at AI research and formal methods roles.
% The full academic CV (complete talk list, full teaching record, GPAs)
% lives in academic-cv.tex.

\documentclass[letterpaper]{article}

\usepackage{geometry}
\usepackage{xcolor}

\usepackage[T1]{fontenc}
\usepackage[sc,osf]{mathpazo}

% Set your name here
\def\name{Shaun Allison}

% Unfilled items render in red so they cannot ship unnoticed.
% Delete every \todo before pushing -- the GitHub Action publishes this PDF.
\newcommand{\todo}[1]{\textcolor{red}{\textbf{[TODO: #1]}}}

\geometry{
  body={6.5in, 8.7in},
  left=1.0in,
  top=1.15in
}

\pagestyle{myheadings}
\markright{\name}
\thispagestyle{empty}

\usepackage{sectsty}
\sectionfont{\rmfamily\mdseries\Large}
\subsectionfont{\rmfamily\mdseries\itshape\large}

\setlength\parindent{0em}

% Tighter lists, no bullets at the top level
\usepackage{enumitem}
\setlist[itemize,1]{label={}, leftmargin=1.2em, topsep=0.35em, itemsep=0.3em,
                    parsep=0em}
\setlist[itemize,2]{label={--}, leftmargin=1.4em, topsep=0.2em, itemsep=0.15em,
                    parsep=0em}
\setlist[enumerate]{leftmargin=1.7em, topsep=0.35em, itemsep=0.35em, parsep=0em}

% hyperref is loaded last, after geometry/sectsty/enumitem, per convention
\usepackage{hyperref}
\hypersetup{
  colorlinks = true,
  urlcolor = [rgb]{0.10,0.20,0.50},
  linkcolor = [rgb]{0.10,0.20,0.50},
  pdfauthor = {\name},
  pdfkeywords = {formal verification, Lean 4, Mathlib, mathematical logic,
                 descriptive set theory, machine learning},
  pdftitle = {\name: Curriculum Vitae},
  pdfsubject = {Curriculum Vitae},
  pdfpagemode = UseNone
}

\begin{document}

{\huge \name}

\vspace{0.16in}

shallison@protonmail.com \quad\textbullet\quad
\url{https://shalliso.github.io/} \quad\textbullet\quad
\url{https://github.com/shalliso}\\
Toronto, Canada\\
US citizen and Canadian permanent resident

\vspace{0.16in}

Mathematical logician and software engineer interested in AI-assisted mathematics, formal verification, and machine learning. Five papers in descriptive set theory, including the \emph{Journal of the European Mathematical Society}, \emph{Advances in Mathematics}, and the \emph{Journal of Mathematical Logic}; Lean~4 formalizations of mathematics; and 14 months as a software engineer at Meta, with two granted patents on machine-learned ranking.

\smallskip
\textbf{Tools:} Lean~4 with Mathlib, Python, \LaTeX{}, git, Pytorch; familiar with numpy,
pandas, matplotlib; earlier industry work in Hack, JavaScript, Haskell,
and C.

\section*{Formal Verification}

\textbf{Knight model formalization in Lean 4 and Mathlib}, 2026\\
\url{https://github.com/shalliso/KnightModel}

\begin{itemize}
  \item Julia Knight's original result was an ad-hoc construction of a
        structure characterizing $\aleph_1$. I generalized this and
        formalized the resulting theory in Lean.
  \item First stage of a project to build higher analogues characterizing
        $\aleph_\alpha$ for every $\alpha < \omega_1$ with AI assistance.
  \item Approximately 5{,}700 lines of Lean, developed with substantial AI
        assistance (Claude) under my direction; the axiomatization, proof
        architecture, and mathematical review are my own.
  \item Presented at the Fields Institute Set Theory Seminar (March 2026), the
        Cornell Logic Seminar (March 2026), and the University of Vienna Logic
        Colloquium (April 2026).
\end{itemize}

\medskip
\textbf{Automatic continuity in Lean 4 and Mathlib}, 2025\\
\url{https://github.com/shalliso/automatic_continuity}

\begin{itemize}
  \item Complete formalization of an automatic continuity theorem:
        every Borel-measurable homomorphism from a Polish group into a
        separable topological group is continuous.
  \item Approximately 300 lines of Lean, written by hand.
\end{itemize}

\section*{Employment}

\begin{itemize}
  \item \textbf{Independent Research}, January 2026 -- present
  \begin{itemize}
    \item Built the Lean 4 formalization described above
    \item Completed revisions on two papers accepted in 2026 (JEMS, JML)
    \item Research visits to Cornell University and the University of Vienna;
          upcoming visit to the Institute for Advanced Study in Mathematics,
          Zhejiang University
    \item Writing a mathematical logic textbook with Spencer Unger
  \end{itemize}

  \item \textbf{Sessional Lecturer I}, University of Toronto,
        September 2025 -- December 2025

  \item \textbf{Postdoctoral Fellow}, University of Toronto,
        August 2022 -- July 2025

  \item \textbf{Postdoctoral Fellow}, Hebrew University of Jerusalem,
        July 2021 -- June 2022

  \item \textbf{Software Engineer}, Meta (then Facebook),
        June 2015 -- August 2016
  \begin{itemize}
    \item Games team: machine-learned targeting of ads and notifications ---
          deciding which users to reach and with what. Two patents granted.
  \end{itemize}

  \item \textbf{Software Engineering Intern}, Meta, Summer 2014
  \item \textbf{Research and Development Intern}, VMware, Summer 2013
  \item \textbf{Intern}, Cisco, Summer 2012
\end{itemize}

\section*{Patents}

\begin{itemize}
  \item \emph{Automatic recipient targeting for notifications} (US10506055B2),
        with L. Foged. Granted. Machine-learned selection of which users
        should receive a given notification.
  \item \emph{Selecting assets} (US11019177B2), with L. Foged. Granted.
        Machine-learned selection of which creative asset to serve.
\end{itemize}

\section*{Publications and Preprints}

\begin{enumerate}
  \item \emph{The class and dynamics of $\alpha$-balanced Polish groups}, with
        Aristotelis Panagiotopoulos.\\
        \textbf{Journal of the European Mathematical Society}, 2026.
        \url{https://doi.org/10.4171/JEMS/1798}

  \item \emph{Countable Borel treeable equivalence relations are classifiable
        by $\ell_1$}.\\
        \textbf{Advances in Mathematics} 479, 2025.
        \url{https://doi.org/10.1016/j.aim.2025.110452}

  \item \emph{Classification strength of Polish groups: Involving $S_\infty$}.\\
        \textbf{Journal of Mathematical Logic}, 2026.
        \url{https://doi.org/10.1142/S0219061326500054}

  \item \emph{Actions of tame abelian product groups}, with Assaf Shani.\\
        \textbf{Journal of Mathematical Logic} 23(3), 2023.
        \url{https://doi.org/10.1142/S0219061322500283}

  \item \emph{Dynamical obstructions to classification by (co)homology and
        other TSI-group invariants}, with Aristotelis Panagiotopoulos.\\
        \textbf{Transactions of the American Mathematical Society} 374, 2021.
        \url{https://doi.org/10.1090/tran/8475}

  \item \emph{Non-Archimedean TSI Polish groups and their potential Borel
        complexity spectrum}.\\
        Preprint, 2020. \url{https://arxiv.org/abs/2010.05085}
\end{enumerate}

\smallskip
\emph{In preparation:} a textbook on mathematical logic, with Spencer Unger.

\section*{Education}

\begin{itemize}
  \item \textbf{Ph.D. Mathematics}, Carnegie Mellon University, 2021
  \begin{itemize}
    \item Dissertation: \emph{Classification of equivalence relations by
          Polish group actions}. Advisor: Clinton Conley
  \end{itemize}
  \item \textbf{B.S. Computer Science}, Carnegie Mellon University, 2015
  \begin{itemize}
    \item Thesis: \emph{A secure human-computable authentication scheme from
          the $k$-Junta problem}. Advisors: Manuel Blum and Jeremiah Blocki
  \end{itemize}
  \item \textbf{B.S. Mathematics} (Discrete Mathematics and Logic),
        Carnegie Mellon University, 2015
\end{itemize}

\section*{Awards}

\begin{itemize}
  \item Putnam Competition, top 500 (2014)
  \item Carnegie Mellon University Honors; College Honors (Computer Science)
  \item Phi Beta Kappa
\end{itemize}

\section*{Selected Talks}

Talks on the Lean formalization are listed under Formal Verification above.
Full list available on request.

\begin{itemize}
  \item \emph{Polish automorphism groups that don't involve $S_\infty$},
        University of Vienna Set Theory Seminar (May 2026)
  \item \emph{The class of CLI Polish groups is coanalytic-complete}, North
        American Descriptive Set Theory Meeting, Brin Mathematics Research
        Center (October 2025)
  \item \emph{A hierarchy of CLI Polish groups and non-reducibility results},
        Caltech Logic Seminar (November 2024)
  \item \emph{A rank determining when a Polish automorphism group has
        less-than-maximal classification strength}, European Set Theory
        Conference (September 2024)
  \item \emph{Polish groups involving $S_\infty$}, Structures: Descriptive Set
        Theory and Dynamics, Institute of Mathematics of the Polish Academy of
        Sciences (August 2023)
  \item \emph{Polish groups which don't involve $S_\infty$}, North American
        Annual Meeting of the Association for Symbolic Logic, UC Irvine
        (March 2023)
  \item \emph{Polish groups with the pinned property}, CMS Winter Meeting,
        Canadian Mathematical Society (December 2020)
  \item \emph{Applications of the Gandy-Harrington topology}, Boise
        Extravaganza in Set Theory, AAAS Pacific Division (June 2019)
\end{itemize}

\section*{Teaching}

\begin{itemize}
  \item Instructor of record for eight courses at the University of Toronto and
        Carnegie Mellon University (2018--2025), from Differential Calculus
        (approx.\ 400 students) to upper-year Topology, Logic, and Field theory.
  \item Supervised an undergraduate independent study in model theory
        (Winter 2023)
  \item Founded a weekly graduate seminar in set theory and model theory at
        Carnegie Mellon University (2019)
\end{itemize}

\bigskip

\begin{center}
  \begin{footnotesize}
    Last updated: \today
  \end{footnotesize}
\end{center}

\end{document}
